% This chapter is the complete working manuscript for the twelfth seminar
% talk. The final section contains supplementary material outside the default
% live route.

\chapter{Elliptic representability I: the smooth-to-derived bridge}
\label[chapter]{chap:Elliptic_Representability_Smooth_Derived_Bridge}

Talk 10 defined the solution stack of a differential moduli problem before
making any analytic choices. Talk 11 began with one elliptic equation near one
solution, chose a Sobolev completion and a finite-dimensional obstruction
space, and produced a Kuranishi model
\[
  \kappa\colon V\longrightarrow C.
\]
Its ordinary zero locus is the local solution set, and its tangent complex is
quasi-isomorphic to the elliptic deformation-obstruction complex.

These two constructions do not yet coincide. The intrinsic solution stack
classifies smooth families over arbitrary test objects. The Kuranishi model
was obtained by applying the Banach implicit function theorem at a chosen
Sobolev regularity. Agreement on ordinary solutions and tangent complexes
does not prove that the model represents the same derived moduli functor.

The purpose of this chapter is to bridge that gap. We localize the intrinsic
stack until it is the derived zero locus of one nonlinear operator on smooth
section spaces. We then use a functorial tower of Sobolev completions to prove
one geometric assertion: the augmented smooth operator is submersive as a
morphism of derived $\Cinfty$-stacks. The Sobolev tower is used inside the
proof, but it is not used to define the moduli object.

This is the technical midpoint of Steffens's representability argument. The
next chapter will take two finite-dimensional pullbacks, construct the local
derived chart, and globalize it.

\section{The theorem and the missing bridge}
\label[section]{sec:Elliptic_Representability_Theorem_Gap}

\subsection{The intrinsic solution stack}
\label[subsection]{sec:Intrinsic_Solution_Stack_Recalled_Twelve}

Let $S$ be a smooth stack and let $M\to S$ be an $S$-family of manifolds. An
elliptic differential moduli problem
\[
  \mathcal E=(X\to M,Y_1\to M,Y_2\to M,P_1,P_2)
\]
is as in \Cref{def:Elliptic_Differential_Moduli_Problem}. The differential
operators induce maps of section stacks
\[
  P_i\colon
  \operatorname{Sect}_{M/S}(Y_i)
  \longrightarrow
  \operatorname{Sect}_{M/S}(X).
\]
Its solution stack is the pullback
\begin{equation}
  \label{eq:Elliptic_Solution_Stack_Twelve}
  \begin{tikzcd}
    \Sol(\mathcal E) \rar \dar \drar[pullback]
      & \operatorname{Sect}_{M/S}(Y_1) \dar["P_1"] \\
    \operatorname{Sect}_{M/S}(Y_2) \rar["P_2"']
      & \operatorname{Sect}_{M/S}(X).
  \end{tikzcd}
\end{equation}
This definition requires neither properness nor ellipticity. It is an
intrinsic construction in the $\infty$-category of derived
$\Cinfty$-stacks.

At a solution $t=(\sigma_1,\sigma_2,\eta)$ over $s\in S$, the relative
tangent complex is
\begin{equation}
  \label{eq:Elliptic_Relative_Tangent_Complex_Twelve}
  \left[
    \begin{matrix}
      \Gamma(M_s,\sigma_1^*T_{Y_{1,s}/M_s})\\
      {}\oplus
      \Gamma(M_s,\sigma_2^*T_{Y_{2,s}/M_s})
    \end{matrix}
    \xrightarrow{D_{\sigma_1}P_1-D_{\sigma_2}P_2}
    \Gamma(M_s,\sigma_X^*T_{X_s/M_s})
  \right]
\end{equation}
in cohomological degrees $0$ and $1$, by
\Cref{thm:Tangent_Complex_Solution_Stack}. Ellipticity says that the displayed
differential is elliptic. If $M_s$ is compact, its kernel and cokernel are
finite-dimensional, although the two chain groups are not.

\subsection{The destination}
\label[subsection]{sec:Relative_Elliptic_Representability_Destination}

We isolate the geometric hypothesis on the domains.

\begin{definition}[Proper family of manifolds]
  \label[definition]{def:Proper_Family_Manifolds_Twelve}
  Let $S$ be a smooth stack. A morphism $M\to S$ is an
  \emph{$S$-family of manifolds} if it is submersive as a morphism of smooth
  stacks. It is a \emph{proper $S$-family of manifolds} if it is also proper.
\end{definition}

For manifolds this is a proper submersion. Ehresmann's theorem then makes it
locally a smooth fiber bundle with compact fiber. Steffens proves the
corresponding local statement over a smooth-stack base
\cite[Proposition~2.1.61]{Steffens_Representability_Elliptic_Moduli}.

\begin{theorem}[Relative elliptic representability]
  \label[theorem]{thm:Relative_Elliptic_Representability}
  Let $S$ be a smooth stack, let $M\to S$ be a proper $S$-family of
  manifolds, and let $\mathcal E$ be an elliptic differential moduli problem
  over $M$. Then
  \[
    \Sol(\mathcal E)\longrightarrow S
  \]
  is representable by quasi-smooth derived $\Cinfty$-schemes locally of
  finite presentation.
\end{theorem}

This is Steffens's Theorem~3.4.3
\cite[Theorem~3.4.3]{Steffens_Representability_Elliptic_Moduli}. We state the
theorem now so that every local construction has a visible destination. Its
proof will be completed in the next chapter.

The conclusion is relative. For every representable map $T\to S$, the base
change
\[
  T\times_S\Sol(\mathcal E)
\]
is represented by a quasi-smooth derived $\Cinfty$-scheme locally of finite
presentation over $T$.

The theorem does not assert compactness, add broken or nodal domains, orient
the determinant complex, or construct a virtual fundamental class. It proves
that the solution functor already defined in
\Cref{eq:Elliptic_Solution_Stack_Twelve} is geometric.

\subsection{Why the Fredholm reduction is not enough}
\label[subsection]{sec:Fredholm_Reduction_Not_Enough_Twelve}

Suppose temporarily that $S$ is a manifold, that $M=S\times N$ with $N$
compact, and that the problem has been placed in vector-bundle coordinates.
It is then governed by a family
\[
  P\colon S\times Q\longrightarrow S\times\Gamma(E;N),
\]
where $Q\subseteq\Gamma(F;N)$ is open. At a solution $(s,\sigma)$, choose a
finite-dimensional space $C$ representing the cokernel of the vertical
linearization and a map
\[
  \iota\colon C\longrightarrow\Gamma(E;N)
\]
with smooth image. The augmented operator is
\begin{equation}
  \label{eq:Augmented_Operator_Informal_Twelve}
  \widetilde P(s',\tau,c)
  =\bigl(s',P_{s'}(\tau)+\iota(c)\bigr).
\end{equation}

Talk 11 applied the implicit function theorem to a Sobolev extension of this
map. That produces an ordinary finite-dimensional augmented solution
manifold, but it leaves three gaps:
\begin{enumerate}
  \item The intrinsic stack uses smooth sections, while the analytic
    construction uses Sobolev sections.
  \item A construction near ordinary points need not classify families over
    derived test objects.
  \item The Kuranishi model depends on auxiliary choices until it is
    identified with an open part of the intrinsic stack.
\end{enumerate}

Smooth sections are the inverse limit of their Sobolev completions in smooth
geometry. The inclusion of smooth stacks into derived $\Cinfty$-stacks need
not preserve such countable limits. Consequently, it would be invalid simply
to take a tower of Sobolev-level derived zero loci and declare its inverse
limit to be the intrinsic solution stack.

Steffens avoids this inverse limit. The tower is used to prove that
\Cref{eq:Augmented_Operator_Informal_Twelve} is a submersion in smooth stacks.
Submersiveness can then be transported into derived geometry, where only
finite pullbacks remain.

\section{Localizing the intrinsic problem}
\label[section]{sec:Localizing_Intrinsic_Elliptic_Problem_Twelve}

\subsection{A fixed compact domain}
\label[subsection]{sec:Fixed_Compact_Domain_Twelve}

Relative representability is local on $S$. Pull back along a smooth atlas,
then shrink around the image of the chosen solution. We may suppose that $S$
is a manifold. By properness and local triviality, we may further identify
\begin{equation}
  \label{eq:Proper_Family_Local_Product_Twelve}
  M\iso S\times N,
\end{equation}
where $N$ is compact.

The two parts of this reduction have different roles. The product description
gives one fixed space of sections and one functorial Sobolev scale throughout
the local family. Compactness makes the Sobolev extensions of elliptic
operators Fredholm.

Fix a solution
\[
  t=(s,\sigma_1,\sigma_2,\eta)
  \in\Sol(\mathcal E).
\]
The path $\eta$ identifies the two output sections and their tangent data. We
write the common output as $\sigma_X$.

\subsection{Section-stack charts}
\label[subsection]{sec:Section_Stack_Charts_Twelve}

A local addition on a submersion $Y\to S\times N$ identifies a neighborhood
of a section $\sigma$ with a neighborhood of the zero section in the vertical
tangent bundle $\sigma^*T_{Y/(S\times N)}$. Since $N$ is compact, this
pointwise construction induces an open chart on the stack of sections.

\begin{proposition}[Local chart on a section stack]
  \label[proposition]{prop:Local_Section_Stack_Chart_Twelve}
  Let $S$ be a manifold, let $N$ be compact, and let
  $Y\to S\times N$ be an $S$-family of submersions. Near a section over
  $s\in S$, there are an open neighborhood $S'\subseteq S$, a vector bundle
  $F\to N$, and an open substack inclusion
  \[
    S'\times Q\longrightarrow
    \operatorname{Sect}_{S'\times N/S'}(Y|_{S'\times N}),
  \]
  where $Q\subseteq\Gamma(F;N)$ is open and contains the zero section.
\end{proposition}

We quote this from Steffens
\cite[Proposition~3.2.22]{Steffens_Representability_Elliptic_Moduli}. The
geometric slogan is that nearby sections are vertical vector fields along the
chosen section. The stack-level assertion says that this description is
compatible with arbitrary parameters and defines an open substack.

Apply the proposition to $Y_1$, $Y_2$, and $X$. Relative jets commute with
the required base changes, by Steffens's Proposition~3.3.9
\cite[Proposition~3.3.9]{Steffens_Representability_Elliptic_Moduli}. The
finite-order operators therefore restrict to maps between the local section
charts.

We will also use the following finite-limit observation.

\begin{lemma}[Open charts pass to pullbacks]
  \label[lemma]{lem:Open_Charts_Pass_To_Pullbacks_Twelve}
  Consider a morphism between two cospans in an $\infty$-category of stacks
  with finite limits:
  \[
    \begin{tikzcd}[column sep=large]
      A \rar \dar[hook] & C \dar[hook] & B \lar \dar[hook] \\
      A' \rar & C' & B' \lar.
    \end{tikzcd}
  \]
  If the three vertical maps are open substack inclusions, then
  \[
    A\times_C B\longrightarrow A'\times_{C'}B'
  \]
  is an open substack inclusion.
\end{lemma}

\begin{proof}
  Open substack inclusions are monomorphisms and are stable under products and
  base change. The induced square
  \[
    \begin{tikzcd}
      A\times_C B \rar \dar \drar[pullback]
        & A\times B \dar \\
      A'\times_{C'}B' \rar
        & A'\times B'
    \end{tikzcd}
  \]
  is a pullback because two points of $A$ and $B$ with equal images in $C'$
  already have equal images in $C$. The result is therefore a base change of
  an open inclusion.
\end{proof}

This is the open-inclusion case of Steffens's Lemma~3.4.4
\cite[Lemma~3.4.4]{Steffens_Representability_Elliptic_Moduli}. Applying it to
the three section charts gives an open substack
\begin{equation}
  \label{eq:Local_Solution_Open_Substack_Twelve}
  \Sol(\mathcal E)_t\longrightarrow\Sol(\mathcal E)
\end{equation}
containing $t$.

\subsection{One nonlinear PDE}
\label[subsection]{sec:One_Nonlinear_PDE_Twelve}

In vector-bundle coordinates, the equality $P_1=P_2$ becomes the zero
equation for their difference. Put
\[
  F=F_1\times_NF_2,
  \qquad
  E=\sigma_X^*T_{X/M}.
\]
The finite-order condition gives an open vector-bundle neighborhood
$W\subseteq J_N^k(F)$ on which the local operator is defined. Let
\[
  Q=\{\tau\in\Gamma(F;N):j^k\tau(N)\subseteq W\}.
\]
Compactness of $N$ makes $Q$ open in the smooth section space. We obtain one
smooth map over $S$,
\begin{equation}
  \label{eq:Local_Nonlinear_PDE_Twelve}
  P\colon S\times Q\longrightarrow S\times\Gamma(E;N),
  \qquad
  (s',\tau)\longmapsto(s',P_{s'}(\tau)).
\end{equation}

The derived zero locus of \Cref{eq:Local_Nonlinear_PDE_Twelve} is precisely
the open stack in \Cref{eq:Local_Solution_Open_Substack_Twelve}. Ellipticity
says that at each solution in this chart the vertical derivative
\begin{equation}
  \label{eq:Vertical_Elliptic_Linearization_Twelve}
  D_\tau P_{s'}\colon
  \Gamma(F;N)\longrightarrow\Gamma(E;N)
\end{equation}
is elliptic.

The local equation has not replaced the intrinsic stack by an analytic
model. It describes an open restriction of that stack before any Sobolev
completion has been chosen.

\section{The Sobolev tower and obstruction augmentation}
\label[section]{sec:Sobolev_Tower_Obstruction_Augmentation_Twelve}

\subsection{The functorial tower}
\label[subsection]{sec:Functorial_Sobolev_Tower_Twelve}

For sufficiently large $l$, the order-$k$ operator in
\Cref{eq:Local_Nonlinear_PDE_Twelve} extends to a smooth map of Banach
manifolds
\begin{equation}
  \label{eq:Sobolev_Level_Operator_Twelve}
  P_l\colon S\times Q_{k+l}
  \longrightarrow S\times H^l(E;N),
\end{equation}
where $Q_{k+l}\subseteq H^{k+l}(F;N)$ is the open subset satisfying the
corresponding jet condition. The shift by $k$ records the loss of derivatives
of an order-$k$ operator.

Steffens's functorial Sobolev-scale lemma supplies four properties which are
used in the proof:
\begin{enumerate}
  \item Smooth sections are the inverse limit of the Sobolev spaces.
  \item The transition maps are compact and have dense image.
  \item Open jet conditions determine compatible open subsets $Q_{k+l}$.
  \item Jet prolongation and the nonlinear operator extend compatibly through
    the tower.
\end{enumerate}
The precise statement is Lemma~3.4.5
\cite[Lemma~3.4.5]{Steffens_Representability_Elliptic_Moduli}. In particular,
we have a commuting tower
\[
  \begin{tikzcd}[column sep=small]
    S\times Q \rar \dar["P"']
      & \cdots \rar
      & S\times Q_{k+l+1} \rar \dar["P_{l+1}"']
      & S\times Q_{k+l} \rar \dar["P_l"']
      & \cdots \\
    S\times\Gamma(E;N) \rar
      & \cdots \rar
      & S\times H^{l+1}(E;N) \rar
      & S\times H^l(E;N) \rar
      & \cdots.
  \end{tikzcd}
\]
Both rows are limit diagrams in convenient manifolds and hence in smooth
stacks.

We use convenient manifolds only as a setting in which smooth section spaces
and the displayed limits exist. No general convenient inverse function
theorem enters the proof.

\begin{warning}[Do not take the derived inverse limit]
  \label[warning]{warn:No_Derived_Sobolev_Inverse_Limit_Twelve}
  The two rows above need not remain limit diagrams after passage from smooth
  stacks to derived $\Cinfty$-stacks. In particular, one may not identify the
  intrinsic solution stack with the inverse limit of the Sobolev-level
  derived zero loci.
\end{warning}

The warning also has a geometric explanation. At one fixed Sobolev level,
reparametrization groups generally fail to act smoothly. A fixed completion
is therefore poorly suited to the global family structure of a moduli
problem. The intrinsic stack remains a stack of smooth sections.

\subsection{A compatible obstruction augmentation}
\label[subsection]{sec:Compatible_Obstruction_Augmentation_Twelve}

Fix the solution $(s,\sigma)$. The Sobolev extension of
$D_\sigma P_s$ is Fredholm. Choose a finite-dimensional space
\[
  C\cong\operatorname{coker}(D_\sigma P_s)
\]
and represent it by smooth sections through a map
\[
  \iota\colon C\longrightarrow\Gamma(E;N).
\]
Define
\[
  \widetilde P_l(s',\tau,c)
  =\bigl(s',P_{l,s'}(\tau)+\iota(c)\bigr).
\]

At $(s,\sigma,0)$, the vertical derivative of $\widetilde P_l$ is
surjective. Openness of the Fredholm locus and upper semicontinuity of
cokernel dimension give compatible neighborhoods
\[
  \widetilde U_l\subseteq S\times Q_{k+l}\times C
\]
on which every $\widetilde P_l$ is a Fredholm submersion. Their smooth cone
point is an open neighborhood
\[
  \widetilde U\subseteq S\times Q\times C.
\]
It carries the smooth augmented operator
\begin{equation}
  \label{eq:Augmented_Smooth_Operator_Twelve}
  \widetilde P\colon
  \widetilde U\longrightarrow S\times\Gamma(E;N),
  \qquad
  (s',\tau,c)\longmapsto
  \bigl(s',P_{s'}(\tau)+\iota(c)\bigr).
\end{equation}

The Banach implicit function theorem applies at each Sobolev level. The
remaining question is whether the smooth map
\Cref{eq:Augmented_Smooth_Operator_Twelve} is submersive in the sense needed
for derived pullbacks.

\section{Testing submersiveness on smooth families}
\label[section]{sec:Testing_Submersiveness_Smooth_Families_Twelve}

\subsection{The categorical criterion}
\label[subsection]{sec:Smooth_Derived_Submersion_Criterion_Twelve}

The comparison we need is a consequence of Steffens's
Proposition~2.2.18(3) and Corollary~2.2.19.

\begin{proposition}[Smooth-to-derived submersion criterion]
  \label[proposition]{prop:Smooth_To_Derived_Submersion_Criterion_Twelve}
  Let $f\colon\mathcal X\to\mathcal Y$ be a morphism of smooth stacks. Suppose
  that for every manifold $Z$ and every map $Z\to\mathcal Y$, the pullback
  formed in smooth stacks
  \[
    Z\times_{\mathcal Y}\mathcal X\longrightarrow Z
  \]
  is represented locally by a submersion of ordinary manifolds. Then the map
  obtained from $f$ in derived $\Cinfty$-stacks is submersive.
\end{proposition}

The result is not merely the assertion that ordinary submersions remain
submersions. Proposition~2.2.18 supplies the representable testing criterion,
and Corollary~2.2.19 transports the resulting notion from smooth to derived
stacks
\cite[Proposition~2.2.18 and Corollary~2.2.19]{Steffens_Representability_Elliptic_Moduli}.

Geometrically, the criterion retains exactly the property needed for finite
pullbacks. A pullback along a submersion is transverse, so its ordinary smooth
pullback already computes the corresponding derived pullback.

\subsection{Sobolev test pullbacks}
\label[subsection]{sec:Sobolev_Test_Pullbacks_Twelve}

\begin{proposition}[Derived submersiveness of the augmented operator]
  \label[proposition]{prop:Augmented_Operator_Derived_Submersive_Twelve}
  After shrinking $\widetilde U$ around $(s,\sigma,0)$, the map
  \Cref{eq:Augmented_Smooth_Operator_Twelve} is submersive as a morphism of
  derived $\Cinfty$-stacks.
\end{proposition}

\begin{proof}
  Let $Z$ be a manifold and let
  \[
    g\colon Z\longrightarrow S\times\Gamma(E;N)
  \]
  be a map of smooth stacks. Compose $g$ with the map to each Sobolev level
  and form
  \begin{equation}
    \label{eq:Sobolev_Test_Pullback_Twelve}
    Z_l
    =Z\times_{S\times H^l(E;N)}\widetilde U_l.
  \end{equation}
  Since $\widetilde P_l$ is a Fredholm submersion, the Banach implicit
  function theorem shows that $Z_l$ is a finite-dimensional manifold and
  $Z_l\to Z$ is a submersion. The spaces $Z_l$ form an inverse tower.

  We show that this tower is eventually locally constant. First consider
  points. A point of $Z_l$ is represented by $(z,\tau,c)$ satisfying
  \[
    P_l(\tau)+\iota(c)=g(z).
  \]
  The right-hand side and $\iota(c)$ are smooth sections. Nonlinear elliptic
  regularity therefore makes $\tau$ smooth. Thus the tower eventually stops
  acquiring new points locally.

  Next consider tangent spaces relative to $Z$. At a smooth point, the
  vertical tangent space of $Z_l\to Z$ is the kernel of the linearized
  augmented operator
  \[
    D\widetilde P_l\colon
    H^{k+l}(F;N)\oplus C\longrightarrow H^l(E;N).
  \]
  Linear elliptic regularity makes every element of this kernel smooth. Hence
  the transition maps eventually induce isomorphisms on vertical tangent
  spaces. Since both stages are submersive over $Z$, their total tangent
  spaces are then isomorphic as well, so the transition map is a local
  diffeomorphism.

  Point stabilization and tangent stabilization together make the inverse
  tower eventually locally constant. Its limit in smooth stacks is therefore
  represented locally by one of the finite-dimensional manifolds $Z_l$, and
  its map to $Z$ is a submersion. Functoriality of the Sobolev tower identifies
  that limit with
  \[
    Z\times_{S\times\Gamma(E;N)}\widetilde U.
  \]
  The criterion in
  \Cref{prop:Smooth_To_Derived_Submersion_Criterion_Twelve} proves the claim.
\end{proof}

The source compresses one analytic point in this proof. Pointwise elliptic
regularity alone does not give a neighborhood on which one Sobolev stage
works uniformly. Steffens's compatible choice of the neighborhoods
$\widetilde U_l$, together with the regularizing property of Lemma~3.4.6,
supplies the required local uniformity
\cite[Proof of Theorem~3.4.3, pp.~69--70]{Steffens_Representability_Elliptic_Moduli}.

The argument has not taken a limit in derived stacks. It took a limit in
smooth stacks, extracted the finite-limit property of submersiveness, and
transported only that property into derived geometry. This is the
smooth-to-derived bridge.

\subsection{The exact handoff}
\label[subsection]{sec:Exact_Handoff_Derived_Chart_Thirteen}

We now know that
\[
  \widetilde P\colon
  \widetilde U\longrightarrow S\times\Gamma(E;N)
\]
is derived-submersive. Consequently, its derived pullback along the zero
section is represented by the corresponding transverse smooth pullback. That
pullback is an ordinary finite-dimensional manifold $V$.

The original equation is not yet obtained by this first pullback, because
$V$ still carries the obstruction coordinate $c$. The next chapter will
project $V$ to $S\times C$ and impose $c=0$ by a second, finite-dimensional
derived pullback. The first pullback removes the infinite-dimensional
analysis. The second retains the finite-dimensional obstruction theory.

\section{Summary and supplementary materials}
\label[section]{sec:Smooth_Derived_Bridge_Summary_Supplementary}

\subsection{Takeaways}
\label[subsection]{sec:Smooth_Derived_Bridge_Takeaways}

The live route can be compressed into eight statements.
\begin{enumerate}
  \item The solution stack is defined intrinsically before choosing a
    completion or obstruction space.
  \item Properness reduces the domain family locally to $S\times N$ with
    $N$ compact.
  \item Local additions and relative jets reduce an open part of the
    intrinsic stack to the derived zero locus of one smooth family $P$.
  \item A functorial Sobolev tower supplies Banach manifolds without
    redefining the moduli problem at finite regularity.
  \item A finite-dimensional obstruction space makes every Sobolev-level
    augmented operator a Fredholm submersion after shrinking.
  \item Smooth test pullbacks are finite-dimensional submersions at every
    level.
  \item Nonlinear and linear elliptic regularity make their towers eventually
    locally constant.
  \item The testing criterion turns this smooth statement into derived
    submersiveness of the augmented operator.
\end{enumerate}

\subsection{Supplement: the general open-pullback lemma}
\label[subsection]{sec:Supplement_General_Open_Pullback_Twelve}

Steffens's Lemma~3.4.4 is more general than
\Cref{lem:Open_Charts_Pass_To_Pullbacks_Twelve}. It considers a morphism of
cospans
\[
  \begin{tikzcd}[column sep=large]
    X_0 \rar \dar["f_0"'] & X_1 \dar["f_1"] & X_2 \lar \dar["f_2"] \\
    Y_0 \rar & Y_1 & Y_2 \lar
  \end{tikzcd}
\]
and a geometric class of morphisms stable under products and base change. If
the $f_i$ lie in that class and the natural comparison conditions hold, then
the induced map
\[
  X_0\times_{X_1}X_2
  \longrightarrow
  Y_0\times_{Y_1}Y_2
\]
lies in the same class. Open inclusions are the case used in the live proof.

\subsection{Supplement: why finite regularity is not the moduli object}
\label[subsection]{sec:Supplement_Finite_Regularity_Not_Moduli_Twelve}

One might try to avoid the Sobolev tower by defining the moduli problem at one
large Sobolev index. This would make the implicit function theorem immediate,
but it would make the global object depend on an arbitrary analytic choice.
It also behaves badly under reparametrization: diffeomorphism groups of the
domain generally do not act smoothly on a fixed Sobolev completion.

Steffens's Remark~3.4.7 emphasizes the alternative. Define the moduli functor
using smooth families, use Sobolev spaces only to prove local
submersiveness, and return to the intrinsic smooth stack before taking the
derived pullback
\cite[Remark~3.4.7]{Steffens_Representability_Elliptic_Moduli}.

\subsection{Related literature}
\label[subsection]{sec:Smooth_Derived_Bridge_Related_Literature}

The primary source is Steffens's proof of Theorem~3.4.3. Proposition~3.2.22
provides section-stack charts, Proposition~3.3.9 controls relative jets,
Lemmas~3.4.4--3.4.6 provide the open-pullback and Sobolev inputs, and
Proposition~2.2.18 together with Corollary~2.2.19 supplies the categorical
bridge
\cite[Proposition~2.2.18, Corollary~2.2.19, Proposition~3.2.22,
Proposition~3.3.9, and Lemmas~3.4.4--3.4.6]{Steffens_Representability_Elliptic_Moduli}.
Pardon's published survey separates the ordinary regularity calculation from
the categorical extension to derived tests by a different method. That
comparison belongs to the end of the next chapter
\cite[Section~5]{Pardon_Representability_Elliptic_Fredholm}.
